\subsubsection{基本解}\label{subsec:fundamental-solutions}

\begin{Definition}
\label{def:fundamental-solution}
设 $L(D)$ 为常系数微分算子. 那么 $L$ 的\textbf{基本解}是满足方程 $L(D)G=\delta$ 的分布 $G\in\mathcal{D}^{\prime}(\mathbb{R}^m)$.
\end{Definition}

当然, 基本解只在相差一个齐次方程 $L(D)u=0$ 的解的意义下唯一; 选择特定基本解时, 常选取在无穷远处具有"最好"行为的那个. 基本解的意义在于如下事实
\begin{equation}\label{eq:fundamental-solution-convolution-identity}
L(D)(G*f)=(L(D)G)*f=\delta*f=f,
\end{equation}
只要卷积 $G*f$ 有定义. 例如, 若 $f$ 具有紧支集, 则 $G*f$ 是方程 $L(D)u=f$ 的一个解.

对一般常系数算子构造基本解相当复杂, 这里将讨论限于一些重要例子.

\begin{Example}{常微分方程}
\label{ex:ode-fundamental-solution}
寻求常微分方程的一个解
\begin{equation}\label{eq:ode-fundamental-solution-equation}
a_nG^{(n)}(x)+\cdots+a_0G(x)=\delta(x).
\end{equation}
对正的和负的 $x$, $G$ 都必须与齐次常微分方程的某个解一致. 即若 $u_1(x)$, $\ldots$, $u_n(x)$ 为齐次常微分方程的一组完备解, 则必须有
\begin{equation}\label{eq:ode-fundamental-solution-piecewise-basis}
G(x)=\begin{cases}\alpha_1u_1(x)+\cdots+\alpha_nu_n(x)&x>0\\\beta_1u_1(x)+\cdots+\beta_nu_n(x)&x<0.\end{cases}
\end{equation}
现在, 要满足 \eqref{eq:ode-fundamental-solution-equation}, 可以要求 $G$ 直到 $(n-2)$ 阶的所有导数都在 0 处连续, 但 $(n-1)$ 阶导数具有大小为 $1/a_n$ 的跳跃. 令 $\gamma_i=\alpha_i-\beta_i$, 得到方程组
\begin{equation}\label{eq:ode-fundamental-solution-jump-system}
\begin{aligned}
\gamma_1u_1(0)+\cdots+\gamma_nu_n(0)&=0,\\
\gamma_1u_1'(0)+\cdots+\gamma_nu_n'(0)&=0,\\
&\vdots\\
\gamma_1u_1^{(n-2)}(0)+\cdots+\gamma_nu_n^{(n-2)}(0)&=0,\\
\gamma_1u_1^{(n-1)}(0)+\cdots+\gamma_nu_n^{(n-1)}(0)&=\dfrac1{a_n}.
\end{aligned}
\end{equation}
这一方程组的行列式是完备解组 $u_i$ 的 Wronskian 行列式, 因而非零.
\end{Example}

\Needspace{4\baselineskip}
\begin{Example}{Laplace 方程}
\label{ex:laplace-fundamental-solution}
现在寻求方程
\begin{equation}\label{eq:laplace-fundamental-solution-equation}
\Delta G(\bm{x})=\delta(\bm{x})
\end{equation}
在 $\mathbb{R}^m$ 上的一个解. 当然, 这使 $G$ 在除原点以外的地方都是齐次 Laplace 方程的解. 此外, 由于 $\delta$ 径向对称, 自然应寻求径向对称的 $G$. 对径向对称函数, Laplace 方程化为
\begin{equation}\label{eq:laplace-radial-fundamental-ode}
G''(r)+\dfrac{m-1}rG'(r)=0,\quad r>0,
\end{equation}
得到解
\begin{equation}\label{eq:laplace-radial-homogeneous-solutions}
G(r)=\begin{cases}c_1+c_2r^{2-m}&m\geq3\\c_1+c_2\ln r&m=2.\end{cases}
\end{equation}
现在可以通过适当选择 $c_2$ 来满足 \eqref{eq:laplace-fundamental-solution-equation}. 对 $m=3$, 在上一节的 Example~\ref{ex:inverse-radius-laplacian} 中已作了这一计算; 结果为 $c_2=-1/4\pi$. 一般结果以类似方式得到; 可以求出基本解
\begin{equation}\label{eq:laplace-fundamental-solution-all-dimensions}
G(\bm{x})=\begin{cases}-r^{2-m}/(m-2)\Omega_m,&m\geq3\\\ln r/2\pi&m=2.\end{cases}
\end{equation}
这里 $\Omega_m=2\pi^{m/2}/\Gamma(m/2)$ 为单位球面的面积 (参见 Problem~\ref{ex:unit-ball-volume-gamma-function}).
\end{Example}
